\section{Artificial Intelligence}
\label{sec:artificial_intelligence}

\subsection*{Introduction}
Artificial Intelligence (\textbf{AI}) refers to the design of systems that can perform tasks usually associated with human intelligence, such as \textbf{perception}, \textbf{reasoning}, \textbf{learning}, \textbf{decision-making}, and language understanding. In modern systems, AI is not a single method but a collection of approaches that range from symbolic reasoning to data-driven learning.

This section provides a compact overview of AI, its main subfields, and the reasons it has become the umbrella concept behind many recent advances in intelligent systems.

\subsection{Main Approaches in AI}
AI systems are usually described through several complementary approaches:

\begin{itemize}
\item \textbf{Symbolic AI:} Uses explicit rules, logic, and knowledge representation.
\item \textbf{Machine Learning:} Learns patterns from data instead of relying only on hand-crafted rules.
\item \textbf{Deep Learning:} Uses multi-layer neural networks to model complex patterns in images, text, audio, and time series.
\item \textbf{Reinforcement Learning:} Learns by interacting with an environment and optimizing long-term reward.
\end{itemize}

\begin{table}[h]
\centering
\caption{Main AI approaches and their typical use cases}
\begin{tabular}{|l|l|l|l|}
\hline
\textbf{Approach} & \textbf{Idea} & \textbf{Strength} & \textbf{Typical Use} \\ \hline
Symbolic AI & Rules and logic & Interpretability & Expert systems \\ \hline
Machine Learning & Learn from data & Adaptability & Prediction and classification \\ \hline
Deep Learning & Multi-layer neural nets & Representation power & Vision and language \\ \hline
Reinforcement Learning & Reward-driven behavior & Sequential decision making & Control and planning \\ \hline
\end{tabular}
\label{tab:ai_approaches}
\end{table}

\textit{[Placeholder for Figure: Taxonomy of Artificial Intelligence approaches]}

\subsection{AI in Modern Applications}
\begin{flushleft}
AI is now used in many practical domains, including healthcare, education, finance, robotics, and natural language processing. The value of AI comes from its ability to automate complex tasks, detect patterns in large datasets, and support human decision-making.

\begin{table}[h]
\centering
\caption{Examples of AI applications across domains}
\begin{tabular}{|l|l|l|}
\hline
\textbf{Domain} & \textbf{Example Task} & \textbf{Benefit} \\ \hline
Healthcare & Diagnosis support & Faster analysis \\ \hline
Education & Personalized learning & Adaptive content delivery \\ \hline
Industry & Predictive maintenance & Reduced downtime \\ \hline
Finance & Fraud detection & Risk reduction \\ \hline
Robotics & Autonomous control & Safe automation \\ \hline
\end{tabular}
\label{tab:ai_applications}
\end{table}

AI systems also depend strongly on data quality, computational resources, and careful validation. Without these, performance can degrade or become unreliable.
\end{flushleft}

\subsection{Limitations and Considerations}
\begin{flushleft}
Although AI enables powerful automation, it introduces several challenges:

\begin{itemize}
\item \textit{Bias:} Models can reproduce or amplify biases present in the training data.
\item \textit{Interpretability:} Complex models are often difficult to explain.
\item \textit{Data Dependence:} Performance depends on the quality and coverage of training data.
\item \textit{Compute Cost:} Large models may require significant memory and processing power.
\end{itemize}
\end{flushleft}

\subsection{Summary}
\begin{flushleft}
Artificial Intelligence provides the general framework within which machine learning, deep learning, and reinforcement learning operate. The following sections focus on the learning-based methods most relevant to this work, starting with reinforcement learning and then moving to neural network architectures such as CNNs and LSTMs.
\end{flushleft}